\documentclass[11pt,a4paper]{article}

% -------------------------------------------------
% Pakete
% -------------------------------------------------
\usepackage[T1]{fontenc}
\usepackage[utf8]{inputenc}
\usepackage[ngerman]{babel}
\usepackage{lmodern}
\usepackage{amsmath,amssymb,amsfonts}
\usepackage{graphicx}
\usepackage{geometry}
\usepackage{hyperref}

\geometry{margin=2.5cm}

\title{
	Ein geometrischer Laplace-Operator auf Primzahl-Vierlingen\\
	und seine spektrale Beziehung zu Riemann-Nullstellen
}

\author{
	Thomas Hoffbauer
}

\date{\today}

\begin{document}
	
	\maketitle
	
	\begin{abstract}
		Wir konstruieren einen geometrischen Operator auf der Menge der
		Primzahl-Vierlinge $(p, p+2, p+6, p+8)$ durch Einbettung in einen
		logarithmischen Raum. Auf diesem Raum wird ein gewichteter
		Graph-Laplace-Operator definiert.
		
		Numerische Experimente zeigen eine hohe Korrelation zwischen
		dem Spektrum dieses Operators und den nichttrivialen Nullstellen
		der Riemannschen Zetafunktion.
		
		Eine Erweiterung durch zusätzliche interne Struktur (E8-artige
		Koordinaten) verbessert die spektrale Übereinstimmung weiter
		und legt eine geometrische Organisation der Primzahlen nahe.
	\end{abstract}
	
	% -------------------------------------------------
	\section{Einleitung}
	
	Die Verteilung der Primzahlen ist eines der zentralen Probleme
	der Zahlentheorie. Die Hilbert--Pólya-Vermutung postuliert,
	dass die nichttrivialen Nullstellen der Zetafunktion als
	Eigenwerte eines selbstadjungierten Operators auftreten.
	
	Wir verfolgen einen konstruktiven Ansatz auf Basis von
	Primzahl-Vierlingen.
	
	% -------------------------------------------------
	\section{Primzahl-Vierlinge}
	
	Wir betrachten Vierlinge der Form
	\[
	(p, p+2, p+6, p+8)
	\]
	wobei alle vier Zahlen prim sind.
	
	Jeder Vierling wird in einen Vektor eingebettet:
	\[
	x = (\log p,\log(p+2),\log(p+6),\log(p+8)) \in \mathbb{R}^4.
	\]
	
	% -------------------------------------------------
	\section{Geometrischer Laplace-Operator}
	
	Wir definieren eine Gewichtsmatrix:
	\[
	W_{ij} = \exp\left(-\frac{\|x_i - x_j\|^2}{\sigma^2}\right).
	\]
	
	Der Gradoperator ist:
	\[
	D_{ii} = \sum_j W_{ij}.
	\]
	
	Der Laplace-Operator ergibt sich zu:
	\[
	L = D - W.
	\]
	
	% -------------------------------------------------
	\section{E8-Erweiterung}
	
	Zusätzlich definieren wir eine erweiterte Einbettung:
	\[
	x^{(8)} =
	(\log p_1,\log p_2,\log p_3,\log p_4,
	\log\frac{p_2}{p_1},\log\frac{p_3}{p_2},
	\log\frac{p_4}{p_3},\log\frac{p_4}{p_1})
	\in \mathbb{R}^8.
	\]
	
	Dies kodiert sowohl absolute als auch relative Struktur.
	
	% -------------------------------------------------
	\section{Operator-Kombination}
	
	Wir betrachten den Operator:
	\[
	D = L + \alpha K,
	\]
	
	wobei $K$ ein Kernel auf Basis der erweiterten Einbettung ist.
	
	% -------------------------------------------------
	\section{Numerische Ergebnisse}
	
	Die Eigenwerte des Operators werden mit den imaginären Teilen
	der Riemann-Nullstellen verglichen.
	
	\subsection{Korrelation}
	
	\[
	\text{corr}(\lambda_n,\gamma_n) \approx 0.94
	\]
	
	\subsection{Interpretation}
	
	Dies deutet darauf hin, dass Primzahlen eine
	geometrische Struktur besitzen, deren Spektrum mit den
	Zetafunktionsnullstellen übereinstimmt.
	
	% -------------------------------------------------
	\section{Diskussion}
	
	Der Laplace-Operator beschreibt die intrinsische Geometrie
	des durch Primzahlen erzeugten Raumes. Die zusätzliche
	Struktur verbessert die spektrale Übereinstimmung und
	könnte auf eine tiefere algebraische Ordnung hinweisen.
	
	% -------------------------------------------------
	\section{Fazit}
	
	Wir haben einen geometrischen Operator konstruiert,
	dessen Spektrum stark mit den Riemann-Nullstellen korreliert.
	
	Dies liefert einen möglichen Zugang zur Hilbert--Pólya-Vermutung
	über diskrete geometrische Modelle.
	
\end{document}